\documentclass{situationist-preprint}

\title{Flew's Gardener Revisited}
\msid{REL-18-b}
\msfamily{REL-18}
\msdiscipline{Philosophy of Religion}
\mstype{Companion manuscript}
\mskeywords{falsification; Flew; Wisdom; parable of the gardener; truth-aptness; religious language}

\begin{document}
\maketitle

\begin{abstract}
Flew's parable of the invisible gardener is the most influential argument that religious assertions, once qualified to survive every possible observation, cease to assert anything. We argue that the parable shows something real but different. The fully qualified gardener hypothesis still excludes a possible state of affairs, that there is no gardener, and it retains inferential connections that the sceptic recognizes; so it remains truth-apt, as REL-18-a concludes of religious statements generally. What the qualifications cost is evidential exposure, not meaning. The parable is an argument about the weight of evidence, and on that question the sceptic may well be right.
\end{abstract}

\section{The parable}

In Wisdom's original, two people return to a neglected garden and find some plants thriving among the weeds; one says that a gardener must have been tending them, the other denies it, and no observation settles their disagreement \citep{Wisdom1945}. Flew sharpened the story. The believer posits a gardener; watches, an electric fence and bloodhounds detect nothing; the believer replies that the gardener is invisible, intangible, insensible to shocks and without scent. The sceptic asks how such a gardener differs from an imaginary gardener or from none at all. Flew's moral is that an assertion compatible with every conceivable state of affairs asserts nothing, and that theological claims risk dying the death of a thousand qualifications \citep{Flew1955}.

The moral runs two properties together. An assertion has \emph{content} if it excludes some possible state of affairs. It is \emph{exposed} if it excludes some observable state of affairs, so that observation could count against it. Flew's principle is plausible for content and false for exposure.

\begin{proposition}\label{p:content}
The fully qualified gardener hypothesis has content and is truth-apt.
\end{proposition}

\begin{proof}
The hypothesis that there is an invisible, intangible gardener who tends the garden excludes the possible state of affairs in which no one tends it and its order arises from soil, weather and chance. It is built by ordinary composition from expressions with ordinary use, its negation is the sceptic's own thesis, and exactly one of the two is true. By Definition 2.1 and Principle P1 of REL-18-a it is truth-apt.
\end{proof}

The qualifications remove exposure, not content. Each excludes a way in which the gardener might have shown himself to the senses; none removes the claim that someone tends the garden. Lemma 3.2 of REL-18-a shows that unexposed truth-apt sentences are commonplace, and the gardener hypothesis is one more.

\section{What the parable does show}

The hypothesis also keeps inferential connections recognized by both parties. If there is a gardener, the order of the plants has an intentional explanation; if the gardener is good, the order is not accidental; if the gardener is powerful and good, the weeds call for an account. The sceptic draws these inferences too, which is why he can ask what the weeds show. Meaning in use, on the view of REL-18-a (Principle P2), is constituted by such connections, and they survive the qualifications.

\begin{objection}
A difference that makes no difference to experience is no difference. If believer and sceptic expect the same observations, they do not disagree about the world.
\end{objection}

\begin{reply}
They disagree about what explains what they both observe. A difference in explanation without a difference in observation is familiar from rival interpretations of the same physical theory, and no one concludes that the interpretations say nothing. Wisdom himself saw that the disagreement was not experimental and did not infer that it was empty.
\end{reply}

What the sceptic legitimately gains is a point about evidence. Each qualification is made to avoid a disconfirming result, and a hypothesis protected in this way receives no support from the observations that might have tested it. The believer who qualifies the gardener out of reach of the watch and the dogs retains a truth-apt claim and gives up any support the garden could have lent it. This is the respect in which Proposition 3.4 of REL-18-a sorts religious claims: some are partly open to evidence and some are not, and the parable is a vivid case of moving a claim from the first class to the second.

\section{Conclusion}

The death by a thousand qualifications is not a death of meaning. The invisible gardener hypothesis says something that may be true, and the sceptic who denies it says something that may be false. What dies is the claim's exposure to the evidence of the garden, and with it any support that evidence could give.

\begin{thebibliography}{99}\small
\bibitem[Flew(1955)]{Flew1955} Flew, A. (1955). Theology and falsification. In A.~Flew and A.~MacIntyre (eds.), \emph{New Essays in Philosophical Theology}. SCM Press.
\bibitem[Wisdom(1945)]{Wisdom1945} Wisdom, J. (1945). Gods. \emph{Proceedings of the Aristotelian Society} 45: 185--206.
\end{thebibliography}

\end{document}
